% FRAGMENT A05 (drafter A: class model and main results)
% REPLACES: main.tex line 203 ("The auction-entry literature makes the bidder pool endogenous" to "… (\citeproc{ref-CarlinEtAl2026}{Carlin et al. 2026}).")
% CHANGE: three short sentences inserted after the first sentence of the line; they place the model in
%   the selective-entry framework. The rest of the line is verbatim.
% NEW CITATION KEYS (bibitems requested from the citation check): ref-MarmerShneyerovXu2013,
%   ref-GentryLi2014. Keys already in the bibliography: ref-RobertsSweeting2013,
%   ref-LevinSmith1994, ref-Samuelson1985.
% SOURCE: audit/alex_meeting_2026_10_21/memos/Q1_literature.md, Claim 1 and Section 7.
%   The two limits are described by the assumptions of each paper (Levin and Smith: only the
%   common distribution before entry; Samuelson: values known before the entry cost). No
%   "nests both" statement is attributed to Gentry and Li. Samuelson's assumption is known
%   through Levin and Smith (1994, footnote 7); the citation check should confirm it.
% LABELS: none kept, removed, or added.
% NUM PLACEHOLDERS: none.
The auction-entry literature makes the bidder pool endogenous (\citeproc{ref-Samuelson1985}{Samuelson 1985}; \citeproc{ref-McAfeeMcMillan1987}{McAfee and McMillan 1987}; \citeproc{ref-LevinSmith1994}{Levin and Smith 1994}; \citeproc{ref-GentryStroup2019}{Gentry and Stroup 2019}) and shows how selective entry affects the choice of sale procedure (\citeproc{ref-RobertsSweeting2013}{Roberts and Sweeting 2013}). The challenger here faces selective entry in the terms of that literature. Its pre-entry information indexes the distribution of its value, and the entry cost reveals the value (\citeproc{ref-MarmerShneyerovXu2013}{Marmer et al. 2013}; \citeproc{ref-GentryLi2014}{Gentry and Li 2014}; \citeproc{ref-RobertsSweeting2013}{Roberts and Sweeting 2013}). The two limits of this framework are a bidder that knows only the common value distribution before entry (\citeproc{ref-LevinSmith1994}{Levin and Smith 1994}) and a bidder that already knows its value (\citeproc{ref-Samuelson1985}{Samuelson 1985}). Here the pre-entry information is the target's stock price, which is public and moves with the orders of an investor who observes the challenger's quality class. I keep the direct deterrence force of that literature and add a market that operates before preparation. Auction formats also affect bidders' incentives to acquire information (\citeproc{ref-Persico2000}{Persico 2000}); the informed trader here is outside the auction and trades a payoff that differs from the entrant's profit. Recent work studies bidder learning about own values and competitors (\citeproc{ref-PernoudGleyze2026}{Pernoud and Gleyze 2026}), post-auction feedback in security-payment design (\citeproc{ref-LiuBernhardt2022}{Liu and Bernhardt 2022}), and bidder-pool choice with correlated values (\citeproc{ref-CarlinEtAl2026}{Carlin et al. 2026}).
